\subsection{Empirical Estimates of AI Productivity}

Эмпирические исследования не позволяют задать единый коэффициент влияния ИИ на разработку. В эксперименте Peng и соавт. участники, завершившие ограниченную задачу на JavaScript с помощью GitHub Copilot, в среднем затратили на неё на 55.8\% меньше календарного времени, чем завершившие задачу участники контрольной группы \cite{peng2023impact}. Результат относится к одной задаче, ранней версии Copilot и выборке, включавшей только завершивших работу участников; параллельное применение нескольких агентов не изучалось. В полевом рандомизированном исследовании Becker и соавт. доступ к инструментам начала 2025 года, напротив, увеличил скорректированное активное время реализации на 19\% у опытных разработчиков зрелых проектов с открытым исходным кодом \cite{becker2025measuring}. Здесь измерялось активное время, а не календарный срок всего набора работ; кроме того, экспериментальное условие включало неоднородные инструменты и способы их применения. Поэтому оценки относятся к разным участникам, задачам, инструментам, правилам включения и показателям результата и не противоречат друг другу.

Различия в эффекте наблюдаются и вне программирования. При внедрении генеративного помощника в службе поддержки Brynjolfsson и соавт. оценили средний рост пропускной способности, измеренной числом решённых обращений в час, в 15\% \cite{brynjolfsson2025generative}. Эффект зависел от опыта сотрудника и типа обращения, а пропускную способность динамического потока нельзя напрямую преобразовать во время фиксированной программной задачи. В терминах модели оценки $k_i^{(r)}$ сопоставимы только при одинаковых границах задачи, исходной конфигурации и критериях приёмки. Календарное время (elapsed time), активное время (active time), пропускная способность (throughput), успешность (success), качество (quality) и понимание результата (comprehension) остаются различными показателями.

Исследования взаимодействия объясняют часть неоднородности. Vaithilingam и соавт. не обнаружили статистически значимого выигрыша Copilot по времени или доле успешно выполненных трёх коротких задач, но выявили затраты на понимание, проверку, редактирование, отладку и переписывание сгенерированного кода \cite{vaithilingam2022expectation}. Barke и соавт. показали, что программисты переключаются между использованием ИИ для ускорения работы и для исследования возможного решения \cite{barke2023grounded}. Поэтому режим $r$ описывает полный способ постановки, исполнения, проверки и исправления, а $h_i^{(r)}$ включает отложенную проверку и повторную работу.

Переход к агентам программирования (coding agents) особенно ясно требует разделять время и качество. В исследовании Balepur и соавт. агент ускорял выполнение исходной веб-задачи и повышал её точность по сравнению с ограниченным чат-ботом, но ухудшал понимание кода; при последующем расширении без агента устойчивого преимущества по времени или точности не обнаружено \cite{balepur2026impaired}. Отсутствие условия без ИИ не позволяет оценить $k$. Телеметрия миллионов сессий агента GitHub Copilot также показывает длинные цепочки вызовов инструментов, слабый внутренний параллелизм и повторные попытки, но не связывает сессии с фиксированными принятыми задачами или действиями человека \cite{liu2026agenticcoding}. Наконец, METR отмечает, что выбор задач с учётом доступности ИИ и параллельная работа с несколькими агентами смещают оценку эффекта на уровне отдельной задачи \cite{metr2026uplift}. Эти результаты обосновывают раздельное измерение автономных и человеческих интервалов и сравнение полных конфигураций рабочего процесса.

\subsection{Effort Estimation as an External Model Layer}

Оценка трудозатрат решает близкую, но самостоятельную задачу. Boehm, Abts и Chulani систематизируют параметрические, регрессионные, аналоговые, экспертные и комбинированные методы прогнозирования трудозатрат, стоимости и длительности \cite{boehm2000survey}. Они оценивают ещё неизвестные входные данные, но сами по себе не назначают неделимые задачи потокам, не проверяют ресурсную допустимость и не минимизируют makespan фиксированного набора работ. Сопоставление формальных моделей и экспертных оценок не выявляет универсально лучшего подхода \cite{jorgensen2007forecasting}; систематическая карта 304 журнальных публикаций также фиксирует преобладание ретроспективной проверки и недостаточное внимание к неопределённости и перспективной проверке (prospective validation) \cite{jorgensenShepperd2007systematic}. Поэтому наблюдаемое после задачи $k_i^{(r)}$ не является готовым прогнозом. До проекта требуется внешняя локально откалиброванная оценка $\hat{k}_i^{(r)}$ с явно заданной неопределённостью; модель расписания преобразует эти оценки в календарные ограничения.

\subsection{Scaling, Orchestration, and Coding-Agent Architectures}

Для фиксированного объёма работ закон Амдала связывает предельное ускорение с последовательной долей \cite{amdahl1967validity}. Здесь используется только структурная аналогия: последовательным ресурсом служит измеренное участие человека, а не участок машинной программы. Универсальный закон масштабируемости Гантера допускает насыщение и отрицательную отдачу из-за конкуренции за ресурсы (contention) и согласования состояния (coherency) \cite{gunther2008general}; его знаменатель является гипотезой для $C(P)$, а не эмпирически установленным законом для агентов программирования. Рассуждение Брукса о коммуникации и концептуальной целостности аналогично обосновывает интеграционные издержки \cite{brooks1995mythical}, но число связей в человеческой команде нельзя буквально переносить на топологию «один разработчик --- несколько агентов».

В универсальных испытаниях многоагентных систем (multi-agent systems, MAS) Kim и соавт. сравнили 260 конфигураций и обнаружили зависимость успешности от делимости задачи, последовательных зависимостей, базовой модели и топологии \cite{kim2025scaling}. Однако результатами служили точность и успешность, человек не участвовал, а число агентов менялось вместе с архитектурой. Это подтверждает зависимость координационных издержек от структуры, но не оценивает $C(P)$ или makespan рассматриваемой человеко-агентной ячейки.

Системы, работающие с репозиториями, дают более близкие, но всё ещё неполные аналоги. MAGIS разделяет роли планировщика, хранителя репозитория, разработчика и контроля качества и допускает автоматизированные циклы проверки и доработки \cite{tao2024magis}. Его основной результат --- доля решённых задач SWE-bench; четыре роли не означают четыре одновременно работающих потока, а автоматизированный контроль качества не является человеческим ресурсом. CAID, напротив, строит план с учётом зависимостей и действительно исполняет агентов-разработчиков конкурентно в изолированных рабочих деревьях Git с последующим слиянием, тестированием и итоговой проверкой \cite{geng2026caid}. В каждой из шести основных пар «семейство модели --- набор задач» CAID повышал показатель качества, но одновременно увеличивал календарное время выполнения и стоимость. Сравнение конфигураций с 2, 4 и 8 агентами показывало немонотонную зависимость качества от их числа и рост издержек. Следовательно, CAID подтверждает различие между заданным числом агентов, фактической конкурентностью и полезным параллелизмом, но не свидетельствует об ускорении и не содержит человеческих фаз, условия без ИИ или разложения $a_i/h_i$.

Таким образом, следует различать внутреннее параллельное использование инструментов одним агентом, одновременное исполнение независимых задач, число именованных ролей и архитектуру координации. Фактическую конкурентность устанавливают только наблюдаемые перекрывающиеся интервалы при неизменном условии сравнения; показатели успешности и качества не заменяют измерения makespan.

\subsection{Scheduling Theory and Repeated Shared Service}

Математическое ядро уровней M1--M2 стандартно. Для одинаковых параллельных машин используются makespan, суммарный объём работы, длительность самой длинной работы и критический путь \cite{pinedo2008scheduling}; списочное планирование при отношениях предшествования и метод LPT для независимых работ имеют классические гарантии \cite{graham1966bounds,graham1969bounds}. Дизъюнктивные дуги, или дуги ресурсного порядка (resource-ordering arcs), также являются стандартным способом представлять конкуренцию машин \cite{pinedo2008scheduling}. Поэтому статья не заявляет новизной makespan, ориентированный ациклический граф (DAG), критический путь, ресурсно-дополненный граф или общий ресурс.

Задача календарного планирования с ограниченными ресурсами (RCPSP) объединяет DAG отношений предшествования с возобновляемыми ресурсами конечной мощности, а многорежимные постановки допускают альтернативные профили длительности и ресурсов \cite{brucker1999resource}. В таком ядре внимание можно представить ресурсом мощности 1, а фазы --- отдельными операциями. Однако стандартная модель не задаёт коэффициент $k$ для программной задачи и режима, смысл принятого результата, операционное разложение $a_i/h_i$ или внутренние интеграционные издержки и издержки переключения.

Ещё ближе модели параллельных машин с общим обслуживающим устройством (common server). В постановках только с подготовкой обслуживающее устройство последовательно подготавливает работы перед параллельной обработкой \cite{hall2000parallel,kravchenko1997parallel}. В постановке с загрузкой и выгрузкой (loading/unloading) каждая работа использует то же устройство до и после машинной обработки; требование немедленной выгрузки может удерживать завершившую обработку машину до освобождения устройства \cite{elidrissi2023loading}. Это формальный аналог последовательности «постановка человеком --- работа агента --- проверка человеком» и частичный аналог удержания потока, но без DAG и повторного цикла «доработка агентом --- проверка человеком».

Ещё более близкая возвратная постановка (reentrant scheduling) проводит работу через назначенную основную машину, удалённый сервер и затем возвращает её на ту же основную машину \cite{chakhlevitch2009reentrant}. Она воспроизводит чередование «исполнитель --- общее обслуживание --- тот же исполнитель», но содержит лишь одно повторное обращение к серверу, независимые работы и заранее заданный маршрут. Существенно, что во время удалённого обслуживания основная машина освобождается; поэтому эта модель не воспроизводит локальную блокировку M4. Ни модели с общим сервером, ни возвратная постановка сами по себе не содержат коэффициент $k$ для программной работы, человеческую приёмку и интеграционные издержки.

Сопоставление уточняет вклад M4: известные механизмы теории расписаний используются как ядро, а предлагаемое отображение на программный процесс связывает полный цикл «постановка человеком --- работа агента --- проверка человеком --- доработка агентом --- приёмка человеком» с режимом задачи, раздельными длительностями фаз, удержанием назначенного агентного потока и фактом принятия результата. Если человеческая фаза не удерживает поток, $A$ и $H$ следует планировать как разные ресурсы; это соответствует другому правилу использования ресурсов.

\subsection{Human Attention in Supervising Autonomous Workers}

Исследования взаимодействия человека и робота (human--robot interaction, HRI) дают содержательный аналог одного оператора и нескольких автономных исполнителей. Показатель числа роботов на одного оператора (fan-out) связывается с продолжительностью взаимодействия и автономной работы без внимания оператора \cite{olsen2004fanout}; модели многозадачного управления требуют укладывать обслуживание одного робота в автономные интервалы остальных \cite{crandall2005validating}. Распределение внимания с учётом состояния показывает зависимость результата от правила выбора и состояния системы, а не только от средней доли участия \cite{crandall2011computing}. Теория многопоточного познания (threaded cognition) дополнительно предупреждает, что не вся человеческая активность единообразно последовательна: интерференция возникает на отдельных когнитивных ресурсах \cite{salvucci2008threaded}.

Cummings и Mitchell проверили расширенную модель числа аппаратов на одного оператора в имитационном эксперименте с 12 участниками и несколькими беспилотными летательными аппаратами (unmanned aerial vehicles, UAV) \cite{cummings2008predicting}. Учёт ожидания взаимодействия, ожидания в очереди, ожидания восстановления ситуационной осведомлённости и когнитивной переориентации снизил расчётную оценку пропускной способности оператора на 36--67\% в рассмотренных условиях; авторы прямо подчёркивают зависимость результата от задачи и контекста. Эксперимент поддерживает последовательную очередь обслуживания и конечность полезного числа автономных исполнителей только как структурную аналогию. Он не калибрует число агентов программирования, $\gamma(P)$ или makespan программной работы: аппараты были однородными и независимыми, а результатом являлась расчётная оценка пропускной способности, а не срок выполнения принятой программной задачи.

\subsection{Synthesis and Bounded Empirical Gap}

Таблица~\ref{tab:closest-prior-work} показывает, что отличие статьи не сводится к добавлению DAG или общего ресурса. Ближайшие частичные аналоги распределены между несколькими направлениями: RCPSP задаёт отношения предшествования и ограничения мощности ресурсов; модели с общим сервером и возвратным обслуживанием --- повторное обслуживание и разные правила удержания машины; HRI --- очередь и переключение одного оператора; универсальные MAS --- влияние топологии координации; MAGIS --- автоматизированный цикл проверки и доработки; CAID --- фактически конкурентные ветви, слияние и тестирование. Каждый из этих элементов известен отдельно.

В основном корпусе, сформированном по поисковому протоколу с датой завершения 2026-08-09, не обнаружена работа, которая одновременно калибровала бы на программных задачах коэффициент $k$ для задачи и режима, автономные и человеческие интервалы, DAG зависимостей, фактическую конкурентность, повторные обращения к единственному человеку, локальное удержание агентного потока, интеграционные издержки и издержки переключения, а также календарный makespan принятого результата. Это утверждение относится только к просмотренному корпусу и каналам поиска, указанным в протоколе; оно не является универсальным утверждением обо всей литературе.

Ближайшая программная система CAID наблюдает конкурентность, календарное время выполнения и интеграцию, но не включает человека и условие без ИИ. MAGIS содержит автоматизированную проверку и доработку, но не устанавливает фактическую конкурентность или человеческое обслуживание. Модели с общим сервером и возвратным обслуживанием формализуют повторный сервис, но не отображают программный процесс; HRI измеряет ожидание человека, но не программную работу; универсальные MAS оценивают преимущественно успешность и качество, а не makespan человеко-агентной ячейки.

Формальный вклад настоящей работы состоит в совместном представлении перечисленных величин и проверке свойств модели на синтетических сценариях. Статья сама не выполняет эмпирическую калибровку реального рабочего процесса. Сбор наблюдений по принятым программным задачам и перспективная проверка (prospective validation) того, предсказывает ли совместно оценённая M4 календарный makespan лучше более простых уровней, остаются будущей работой.

\begin{landscape}
\begin{table}[!ht]
\centering
\caption{Ближайшие классы постановок и область настоящей модели. Значения «да», «частично» и «нет» показывают наличие свойства в классе постановок, а не эмпирическую эквивалентность или новизну.}
\label{tab:closest-prior-work}
\setlength{\tabcolsep}{2.1pt}
\renewcommand{\arraystretch}{1.18}
\renewcommand{\textsuperscript}[1]{\ensuremath{^{\mathrm{#1}}}}
\scriptsize
\begin{tabular}{>{\raggedright\arraybackslash}p{2.85cm}>{\centering\arraybackslash}p{1.45cm}>{\centering\arraybackslash}p{1.9cm}>{\centering\arraybackslash}p{1.65cm}>{\centering\arraybackslash}p{1.7cm}>{\centering\arraybackslash}p{1.6cm}>{\centering\arraybackslash}p{1.45cm}>{\centering\arraybackslash}p{1.55cm}>{\centering\arraybackslash}p{1.95cm}>{\raggedright\arraybackslash}p{2.6cm}>{\raggedright\arraybackslash}p{4.15cm}}
\toprule
Класс постановок & DAG & \shortstack{Потоки\\исполнителей} & \shortstack{Сервис\\мощности 1} & \shortstack{Повторный\\сервис} & \shortstack{Удержание\\потока} & \shortstack{$k$ задачи\\и режима} & \shortstack{Разделение\\$a_i/h_i$} & \shortstack{Издержки\\координации\\и интеграции} & \shortstack{Основной\\результат} & \shortstack{Калибровка реального\\программного процесса} \\
\midrule
RCPSP и многорежимная RCPSP \cite{brucker1999resource}\textsuperscript{a}
& да & частично & да & частично & нет & нет & нет & частично & makespan, стоимость, допустимость & нет \\

Параллельные машины с общим сервером \cite{hall2000parallel,kravchenko1997parallel,elidrissi2023loading,chakhlevitch2009reentrant}\textsuperscript{b}
& нет & да & да & частично & частично & нет & частично & нет & makespan, алгоритмические оценки & нет \\

HRI: число роботов на оператора и распределение внимания \cite{olsen2004fanout,crandall2005validating,crandall2011computing,cummings2008predicting}\textsuperscript{c}
& нет & да & да & да & частично & нет & частично & частично & пропускная способность, выполнение миссии & нет \\

Универсальные MAS и централизованная координация \cite{kim2025scaling}\textsuperscript{d}
& нет & частично & нет & частично & нет & нет & нет & частично & успешность, качество; иногда стоимость & нет \\

Агенты программирования: MAGIS и CAID \cite{tao2024magis,geng2026caid}\textsuperscript{e}
& частично & частично & нет & частично & нет & нет & нет & частично & успешность, качество; CAID также время и стоимость & нет \\

Настоящая работа\textsuperscript{f}
& да & да & да & да & да & да & да & да & makespan принятого результата & нет \\
\bottomrule
\end{tabular}

\vspace{0.5em}
\begin{minipage}{0.97\linewidth}
\scriptsize
\textsuperscript{a} Возобновляемый ресурс RCPSP не обязан представлять одинаковые агентные потоки; повторное обслуживание и издержки требуют явных операций, а стандартная постановка не вводит автоматическое удержание потока в ожидании другого ресурса.

\textsuperscript{b} Повторный сервис присутствует в вариантах с загрузкой/выгрузкой и возвратом, но не в постановках только с подготовкой. Немедленная выгрузка может удерживать машину, тогда как возвратная модель с удалённым сервером освобождает её; машинная и сервисная длительности не откалиброваны как доли работы агента и человека.

\textsuperscript{c} Автономная система ожидает оператора или теряет результативность, однако это не равнозначно удержанию программного потока. Интервалы взаимодействия и автономности не являются программными $a_i/h_i$; очередь, потеря ситуационной осведомлённости и переориентация не включают интеграцию кода.

\textsuperscript{d} Последовательная взаимозависимость измеряется, но явный DAG задач не строится; число агентов или ролей не доказывает перекрытие интервалов. Централизованный координатор не равнозначен наблюдаемому сервису мощности 1, повторные сообщения --- человеческому циклу, а различия результатов --- калибровке издержек M4.

\textsuperscript{e} CAID строит план зависимостей и показывает перекрывающиеся ветви; для MAGIS и класса в целом эти свойства не установлены. Автоматизированные планировщик, контроль качества и доработка не являются человеческим сервисом; ожидание зависимостей или слияния не тождественно локальной блокировке. CAID наблюдает слияние, тестирование, время и стоимость, но не калибрует человеко-агентный процесс. Оценивание на программных наборах задач не является эмпирической калибровкой реального рабочего процесса с наблюдаемыми человеческими фазами.

\textsuperscript{f} В статье используются синтетические сценарии; эмпирическая калибровка и перспективная проверка остаются будущей работой.

\medskip
Статья не заявляет новизной DAG, makespan, общий ресурс, повторный сервис, возврат к исполнителю, удержание ресурса, дуги ресурсного порядка или число автономных исполнителей на одного оператора.
\end{minipage}
\end{table}
\end{landscape}

\FloatBarrier
